\section{结论}
本文给出了 Long CoT 学习的机制解释：三类“键”——自我反思、深层推理、自我探索——在推理过程中逐渐形成类似分子的结构。我们基于行为转移分布提出“语义异构体”框架，从而洞察 Long CoT 学习的稳定性与失效模式。在此基础上，我们提出 \method，利用分布迁移图构建稳健的 Long CoT 结构，显著提升推理性能并增强强化学习的稳定性。

\newpage
\section*{局限性}
尽管我们在多个推理基准上取得了不错的表现，仍存在若干局限：首先，受成本与规模限制，我们的分析仅覆盖有限的教师模型与学生骨干，因此观测到的 Long CoT 统计模式可能偏向某些架构或训练配方。其次，我们聚焦线下蒸馏与监督微调，尚未验证该方法在带有 RL 回馈的在线或交互式场景中的扩展性。再次，我们只能近似可视化信息空间与语义空间中推断出的键几何特征，要精准刻画普适的长链大分子结构仍是重要的未来方向。最后，我们的行为分析依赖自动标注数据，尽管进行过初步的鲁棒性检查，仍不可避免地存在噪声或偏差。
